\documentclass[10pt]{article}
\usepackage[letterpaper,margin=0.7in,top=0.6in,bottom=0.6in]{geometry}
\usepackage[T1]{fontenc}
\usepackage[utf8]{inputenc}
\usepackage{charter}
\usepackage[scaled=0.9]{helvet}
\usepackage{microtype}
\usepackage{xcolor}
\usepackage{titlesec}
\usepackage{enumitem}
\usepackage{tabularx}
\usepackage{etoolbox}
\usepackage[hidelinks]{hyperref}

\definecolor{accent}{HTML}{0F6A73}
\definecolor{ink3}{HTML}{6B7177}
\hypersetup{colorlinks=true,urlcolor=accent,linkcolor=accent}
\pagestyle{empty}
\setlength{\parindent}{0pt}
\linespread{1.06}

\titleformat{\section}{\sffamily\bfseries\small\color{accent}}{}{0pt}{\MakeUppercase}[\vspace{-6pt}{\color{accent!35}\rule{\linewidth}{0.5pt}}]
\titlespacing*{\section}{0pt}{12pt}{6pt}

% \entry{when}{title}{where}{description}
\newcommand{\entry}[4]{%
  \noindent\begin{tabularx}{\linewidth}{@{}p{0.95in}X@{}}
  {\leavevmode\sffamily\small\color{ink3}#1} & \textbf{#2}\ifblank{#3}{}{\ \textperiodcentered\ #3}
  \ifblank{#4}{}{\newline{\small #4}}
  \end{tabularx}\par\vspace{5pt}}

% \pub{venue}{badge}{title}{authors}{summary}{links}
\newcommand{\pub}[6]{%
  \noindent\begin{tabularx}{\linewidth}{@{}p{0.95in}X@{}}
  \raggedright{\sffamily\small\bfseries #1}\ifblank{#2}{}{\newline{\sffamily\footnotesize\color{ink3}#2}}
  & \textbf{#3}\newline{\small #4}\newline{\small #5}\ifblank{#6}{}{\newline{\footnotesize #6}}
  \end{tabularx}\par\vspace{6pt}}
\newcommand{\me}{\underline{Boaz Meivar}}
\newcommand{\qone}{{\color{accent}\textbf{Q1}}}

\begin{document}

{\sffamily\bfseries\LARGE Boaz Meivar}\par\vspace{3pt}
{\small PhD student, computer vision, Tel Aviv University \textperiodcentered\ advised by Prof.\ Shai Avidan \textperiodcentered\ Tel Aviv, Israel}\par\vspace{2pt}
{\small\href{mailto:bmeivar@gmail.com}{bmeivar@gmail.com} \quad
\href{https://bmeivar.github.io}{bmeivar.github.io} \quad
\href{https://github.com/bmeivar}{github.com/bmeivar} \quad
\href{https://www.linkedin.com/in/boaz-meivar}{linkedin.com/in/boaz-meivar} \quad
\href{https://www.semanticscholar.org/author/2148726410}{Semantic Scholar}}

\section{Research}
I study when and why visual representations break: how object retrieval degrades with object scale and context, how to evaluate retrieval at web scale at a fraction of the cost, and whether vision-language models use the evidence in front of them.

\vspace{4pt}{\sffamily\small\bfseries Key results}
\begin{itemize}[leftmargin=1.2em,itemsep=1pt,topsep=2pt]
\item \textbf{State of the art on ILIAS-100M} instance retrieval (42.0 mAP@1000), training-free and without re-indexing the gallery.
\item \textbf{Web-scale retrieval evaluation with 20$\times$ less indexing}, accurate to about 0.01 mAP.
\item \textbf{Ill-posed tasks as a new lens on vision-language models:} when no single cue determines the answer, counterfactuals reveal which evidence a model actually uses.
\end{itemize}

\section{Education}
\entry{2024--now}{PhD, Electrical Engineering}{Tel Aviv University}{Research: AI and computer vision \textperiodcentered\ advisor: Prof.\ Shai Avidan \textperiodcentered\ expected 2028}
\entry{2020--2022}{MSc, Aerospace Engineering}{Technion -- Israel Institute of Technology}{Advisor: Prof.\ Moti Karpel \textperiodcentered\ thesis: data-driven aeroelastic flutter precursor}
\entry{2008--2013}{BSc, Aerospace Engineering}{Technion -- Israel Institute of Technology}{Dean's List}

\section{Experience}
\entry{2022--now}{Deep Learning Algorithm Researcher}{computer vision stealth startup}{Originated and led research on retrieving arbitrary objects in massive geospatial datasets from as little as one example, and built an MVP, now a cornerstone product. Combines active learning, self-supervised foundation models and gradient-boosted trees. Patent pending.}
\entry{2020--2022}{Computer Vision Algorithm Developer}{aerospace industry}{Maintained and optimized a dim-target detection algorithm, wrote simulations in MATLAB, and ran operational analyses to improve its performance. Started and taught a 24-hour Python workshop for the engineering team.}
\entry{2014--2020}{Dynamics and Loads Engineer}{aerospace industry}{Multi-body dynamics simulations, calibrated with genetic optimization algorithms.}

\section{Publications}
{\small\color{ink3}* equal contribution}\par\vspace{4pt}
{\sffamily\small\bfseries Computer vision and machine learning}\par\vspace{4pt}
\pub{WACV 2027}{accepted}{Object-Only Is Not Context-Free: A Counterfactual Diagnostic for Context Removal in Object Retrieval}
{\me, Shai Avidan}
{KAP, the Keyhole Aperture Probe, measures how much an object's embedding depends on its surroundings, by revealing more of the scene around the object while the object region used for pooling stays fixed.}{}
\pub{arXiv 2026}{under review}{Beyond Resolution: Object-to-Image Ratio Mismatch in Instance Retrieval}
{\me, Ofir Kedem, Amit Edenzon, Gal Chechik, Shai Avidan}
{Object-to-image ratio mismatch, not resolution, explains over 80\% of cross-distance retrieval loss for 9 of 12 backbones. Query-side scale augmentation reaches 29.2 mAP@1000 on ILIAS-100M, and adding an OWLv2 crop reranker gives a state-of-the-art 42.0, all training-free and without re-indexing the gallery.}
{\href{https://arxiv.org/abs/2610.11489}{arXiv:2610.11489}}
\pub{arXiv 2026}{under review}{Web Scale on a Budget: Tail-Aware Retrieval Evaluation}
{\me, Shai Avidan}
{From a 5\% sample, a tail-aware (EVT) estimate lands within about 0.01 of full-scale mAP (MAE 0.005) across five backbones, validated at 100M images on ILIAS and INSTRE and on the full 759k-distractor GLDv2 pool.}
{\href{https://arxiv.org/abs/2610.17986}{arXiv:2610.17986}}
\pub{arXiv 2026}{under review}{Ill-Posed by Design: Probing Evidence Use in VLMs}
{\me, Shaked Perek, Shani Shvartzman, Eli Schwartz, Shai Avidan}
{Monocular metric object size is ill-posed, so counterfactuals separate unused evidence from redundant evidence. Across 12 open VLMs (3B to 397B) and six evidence channels, object identity is the strongest cue and global scene geometry is largely unused.}
{\href{https://arxiv.org/abs/2606.24335}{arXiv:2606.24335} \textperiodcentered\ \href{https://github.com/bmeivar/MetricVQA}{code}}
\pub{WACV 2026}{}{Decoupling Shape and Texture in SAM-2 via Controlled Texture Replacement}
{Inbal Cohen*, \me*, Peihan Tu, Shai Avidan, Gal Oren}
{Fine-tuning with incrementally stronger texture replacement gives TextureSAM, which segments better where texture defines the boundary.}
{\href{https://doi.org/10.1109/WACV61042.2026.00397}{paper} \textperiodcentered\ \href{https://arxiv.org/abs/2505.16540}{arXiv:2505.16540} \textperiodcentered\ \href{https://github.com/Scientific-Computing-Lab/TextureSAM}{code}}

\vspace{2pt}{\sffamily\small\bfseries Machine learning for science and engineering}\par\vspace{4pt}
\pub{Nature Sci.\ Rep.\ 2026}{\qone}{Measurement-aware learning for reliable grain-boundary analysis in quantitative metallography}
{\me, Inbal Cohen, Matan Rusanovsky, Francis Quintal Lauzon, Julien Robitaille, Ofer Be'eri, Shai Avidan, Gal Oren}
{Removing global context sharply degrades detection of weak grain boundaries; MLOGRAPHY++ is validated on the grain-size measurement itself, not only on pixel accuracy.}
{\href{https://doi.org/10.1038/s41598-026-58714-5}{doi:10.1038/s41598-026-58714-5}}
\pub{Annals BME 2026}{}{Simultaneous 3D Quantification of Fluid--Structure Interaction in a Patient-Averaged Left Ventricle: A Lagrangian Experimental Framework}
{Omer Hadar*, \me*, Shai Avidan, Alex Liberzon}
{SAM-2 segmentation with 3D particle tracking gives co-registered, time-resolved 3D wall motion and intraventricular flow in a compliant heart model.}
{\href{https://doi.org/10.1007/s10439-026-04293-5}{doi:10.1007/s10439-026-04293-5}}
\pub{AIAA J.\ 2025}{\qone}{Data-Driven Approach to the Development of an Aeroelastic Flutter Precursor}
{\me, Moti Karpel}
{A flutter early-warning system built end to end with essentially no prior knowledge of the aircraft; thresholds set on a wind-tunnel wing carried over to a flight test, warning up to 6\,s before critical flutter onset.}
{\href{https://doi.org/10.2514/1.J064327}{doi:10.2514/1.J064327} \textperiodcentered\ \href{https://github.com/bmeivar/flutter}{code}}

\section{Skills}
\begin{tabularx}{\linewidth}{@{}p{0.95in}X@{}}
{\leavevmode\sffamily\small\color{ink3}Methods} & \small Self-supervised learning, few-shot retrieval, active learning, gradient-boosted trees, LoRA fine-tuning\\[2pt]
{\leavevmode\sffamily\small\color{ink3}Models} & \small DINOv2, DINOv3, SigLIP2, SAM / SAM-2, OWLv2, open VLMs (Qwen-VL and InternVL families)\\[2pt]
{\leavevmode\sffamily\small\color{ink3}Engineering} & \small Python, PyTorch, Hugging Face Transformers, FAISS; multi-GPU feature extraction and memory-mapped descriptor stores for 100M-image search\\[2pt]
{\leavevmode\sffamily\small\color{ink3}Also} & \small Time-series analysis and anomaly detection, physics simulation and optimization; MATLAB\\[2pt]
{\leavevmode\sffamily\small\color{ink3}Languages} & \small Hebrew, English, French
\end{tabularx}

\end{document}
